\section{Dữ liệu, công cụ và quy trình}
\label{sec:data}

\subsection{Bối cảnh và phạm vi dữ liệu}

Sông Colorado là một trong những hệ thống sông quan trọng nhất của miền Tây Hoa Kỳ.
Lưu lượng dòng chảy tại một trạm quan trắc phản ánh tổng hợp nhiều quá trình: lượng
mưa và tuyết rơi trên lưu vực, tốc độ tan tuyết theo mùa, đặc điểm địa hình, lớp phủ
thực vật và cả hoạt động điều tiết của hồ chứa. Chuỗi lưu lượng ngày vì vậy thường
có tính dai dẳng cao: giá trị hôm nay và những ngày gần nhất có tương quan mạnh với
giá trị ngày mai. Tuy nhiên, chuỗi cũng có thể thay đổi mức nền theo năm do điều kiện
khí hậu và vận hành hệ thống, tạo ra xu hướng dài hạn và các giai đoạn lưu lượng cao
bất thường.

Báo cáo sử dụng tệp \texttt{colorado\_river\_flow.csv} đi kèm bài tập. Tệp có hai cột:
\texttt{date} (ngày quan trắc) và \texttt{discharge\_m3s} (lưu lượng, đơn vị
m\textsuperscript{3}/s). Khoảng thời gian từ \DateMin{} đến \DateMax{}, gồm
\DatasetN{} ngày. Bộ dữ liệu không kèm mô tả chi tiết về trạm đo, đơn vị vận hành
hay quy trình hiệu chỉnh số liệu; vì vậy báo cáo chỉ phân tích cột lưu lượng được
cung cấp và không suy diễn thêm về nguồn gốc hoặc chất lượng đo.

Mục tiêu của phần dữ liệu là trả lời ba câu hỏi kỹ thuật trước khi mô hình hóa:
(1) chuỗi có đủ dài và đủ liên tục để phân tích mùa vụ và dự báo hay không;
(2) có ngày thiếu, ngày trùng, giá trị null hoặc giá trị bất thường về miền xác định
hay không; (3) cách chia Train--Validation--Test có bảo toàn thứ tự thời gian và
tránh rò rỉ thông tin hay không. Các câu hỏi này được kiểm tra tự động trong pipeline
và kết quả được lưu thành tệp JSON để báo cáo có thể tái lập.

\subsection{Cấu trúc và cách đọc dữ liệu}

Bảng~\ref{tab:schema} mô tả hai cột của dữ liệu gốc. Cột ngày được chuyển sang kiểu
\texttt{datetime} và sắp xếp tăng dần; cột lưu lượng được giữ ở kiểu số thực. Sau khi
kiểm tra, dữ liệu được ghi lại thành tệp đã xử lý để các bước sau không phải đọc lại
và diễn giải lại dữ liệu thô.

\begin{table}[H]
\centering
\small
\begin{tabularx}{\linewidth}{@{}llX@{}}
\toprule
\textbf{Cột} & \textbf{Kiểu} & \textbf{Ý nghĩa và cách xử lý} \\
\midrule
\texttt{date} & ngày & Ngày quan trắc, định dạng ISO \texttt{YYYY-MM-DD}; dùng làm trục thời gian và khóa duy nhất. \\
\texttt{discharge\_m3s} & số thực & Lưu lượng dòng chảy, đơn vị m\textsuperscript{3}/s; giữ nguyên giá trị gốc, không làm tròn khi tính toán. \\
\bottomrule
\end{tabularx}
\caption{Mô tả các cột dữ liệu gốc và cách xử lý.}
\label{tab:schema}
\end{table}

Vì dữ liệu chỉ có một chuỗi lưu lượng, các phân tích mô tả được thực hiện trên toàn
bộ 14 năm. Bài toán dự báo sử dụng cùng chuỗi nhưng chia theo thời gian: Train
2010--2019, Validation 2020--2021 và Test 2022--2023. Cách chia này mô phỏng tình
huống thực tế: khi dự báo một ngày trong tương lai, mô hình chỉ được phép dùng thông
tin có trước thời điểm đó. Không có phép trộn ngẫu nhiên giữa các năm, vì trộn ngẫu
nhiên sẽ đưa thông tin tương lai vào tập huấn luyện và làm sai lệch đánh giá.

\subsection{Công cụ và môi trường}

Pipeline được viết bằng Python 3.12, sử dụng \texttt{pandas} \cite{mckinney2010} và
\texttt{numpy} \cite{harris2020} cho xử lý dữ liệu, \texttt{statsmodels}
\cite{seabold2010} cho phân rã STL và tự tương quan, \texttt{lightgbm}
\cite{ke2017} cho mô hình cây tăng cường gradient, \texttt{matplotlib}
\cite{hunter2007} cho hình tĩnh và \texttt{pytest} cho kiểm thử. Các tham số chính được đặt trong
\texttt{code/configs/default.yaml}; phiên bản thư viện được ghi trong
\texttt{code/requirements.lock.txt}. Website kết quả dùng React, TypeScript, Vite và
ECharts, được xây dựng thành các tệp tĩnh và triển khai bằng GitHub Pages.

Báo cáo được biên dịch bằng XeLaTeX thông qua Tectonic. Bảng và hình trong báo cáo
không được chép tay từ màn hình: chúng được sinh từ kết quả pipeline, sau đó LaTeX
chỉ việc chèn nội dung đã sinh. Cách làm này giúp giảm sai lệch giữa số liệu trong
báo cáo, số liệu trên website và số liệu trong mã nguồn. Mọi lần chạy lại pipeline
đều tạo một \texttt{run\_id}; báo cáo dùng run \RunId{} để đối chiếu.

\subsection{Quy trình tái lập}

Pipeline gồm bảy bước có thứ tự:

\begin{enumerate}[leftmargin=1.6em]
  \item \textbf{Audit:} đọc dữ liệu thô, kiểm tra ngày thiếu/trùng, null, giá trị
  âm, giá trị bằng 0 và tính liên tục của lịch.
  \item \textbf{EDA:} tính thống kê mô tả, thống kê theo tháng/năm, phân rã STL và
  tự tương quan.
  \item \textbf{Features:} tạo lag, rolling window, đặc trưng lịch và biến mục tiêu
  cho từng horizon.
  \item \textbf{Train:} huấn luyện LightGBM trên Train, chọn tham số bằng Validation.
  \item \textbf{Freeze:} khóa cấu hình, tham số và split để tránh thay đổi sau khi
  đã xem Test.
  \item \textbf{Evaluate:} dự báo trên Test, tính RMSE, MAE, NSE theo từng horizon,
  theo năm và theo mức lưu lượng.
  \item \textbf{Export:} sinh bảng, hình, số liệu LaTeX, manifest và dữ liệu cho
  website.
\end{enumerate}

Mỗi bước ghi kết quả trung gian vào \texttt{code/results}. Nếu một bước bị lỗi, các
bước sau không được chạy tiếp. Cách tổ chức này quan trọng với báo cáo vì nó buộc
mọi con số phải có nguồn gốc rõ ràng: từ dữ liệu thô, qua cấu hình, đến tệp kết quả
và cuối cùng là bảng/hình trong LaTeX.

\subsection{Tiêu chí đánh giá dự báo}

Ba thước đo được dùng là RMSE, MAE và NSE; NSE được đề xuất trong
\cite{nash1970}. Với $n$ cặp giá trị thực $y_i$ và dự báo $\hat{y}_i$:

\begin{equation}
\mathrm{RMSE} = \sqrt{\frac{1}{n}\sum_{i=1}^{n}(y_i-\hat{y}_i)^2},
\qquad
\mathrm{MAE} = \frac{1}{n}\sum_{i=1}^{n}|y_i-\hat{y}_i|.
\label{eq:rmse-mae}
\end{equation}

\begin{equation}
\mathrm{NSE} = 1-\frac{\sum_{i=1}^{n}(y_i-\hat{y}_i)^2}
{\sum_{i=1}^{n}(y_i-\bar{y})^2}.
\label{eq:nse}
\end{equation}

RMSE phạt mạnh các sai số lớn và có cùng đơn vị với lưu lượng; MAE dễ diễn giải hơn
vì là sai số tuyệt đối trung bình; NSE so sánh mô hình với một dự báo “trung bình”
của chính tập đánh giá. NSE bằng 1 là hoàn hảo, bằng 0 là không tốt hơn trung bình,
âm là kém hơn trung bình. Vì các horizon có độ khó khác nhau, báo cáo trình bày cả
trung bình theo horizon (mean) và giá trị gộp trên toàn bộ cặp dự báo (pooled). Hai
cách tính này không thay thế nhau: mean cho biết mô hình tốt đến đâu ở từng horizon,
pooled cho biết sai số tổng hợp trên toàn bộ tập Test.
